\section{Main results}\label{sec:results}

\begin{theorem}[fixed-side short-lifetime law]\label{thm:main}
For the normalized side-24 field \eqref{eq:K24}, the finite nonessential superlevel $H_0$ first-moment lifetime density satisfies
\begin{equation}\tag{2.1}\label{eq:main}
\nu_{3,24}(\ell)=c_{3,24}\,\ell^{-1/3}\bigl(1+O(\ell^{1/3})\bigr),\qquad \ell\downarrow0,
\end{equation}
equivalently $\nu_{3,24}(\ell)=c_{3,24}\ell^{-1/3}+O(1)$, where $0<c_{3,24}<\infty$ is given by \eqref{eq:c} below.

Fix $t\in S^2$ and an orthonormal transverse frame $n_1,n_2$. Define
\begin{equation}\tag{2.2}
u_t=-\frac{\partial_t^3f_{24}}{12},\qquad g_t=(\partial_tf_{24},\partial_{n_1}f_{24},\partial_{n_2}f_{24}),
\end{equation}
\begin{equation}\tag{2.3}
h_t=(\partial_{tt}f_{24},\partial_{tn_1}f_{24},\partial_{tn_2}f_{24}),
\end{equation}
and
\begin{equation}\tag{2.4}
q_t=(\partial_{n_1n_1}f_{24},\partial_{n_1n_2}f_{24},\partial_{n_2n_2}f_{24}),\qquad
Q_t(q)=\begin{pmatrix}q_1&q_2\\ q_2&q_3\end{pmatrix}.
\end{equation}
Put
\begin{equation}\tag{2.5}
\Sigma_g(t)=\Cov(g_t),\qquad \Sigma_h(t)=\Cov(h_t),
\end{equation}
\begin{equation}\tag{2.6}
\sigma_t^2=\Var(u_t)-\Cov(u_t,g_t)\,\Sigma_g(t)^{-1}\Cov(g_t,u_t),
\end{equation}
\begin{equation}\tag{2.7}
\Gamma_t=\Cov(q_t)-\Cov(q_t,h_t)\,\Sigma_h(t)^{-1}\Cov(h_t,q_t),
\end{equation}
and
\begin{equation}\tag{2.8}\label{eq:Dt}
D_t=\E\Bigl[(\det Q_t(q_t))^2\,\ind_{\{Q_t(q_t)<0\}}\ \Big|\ h_t=0\Bigr].
\end{equation}
Then
\begin{equation}\tag{2.9}\label{eq:c}
c_{3,24}=B_3\int_{S^2}\frac{(\sigma_t^2)^{2/3}D_t}{\sqrt{\det\Sigma_g(t)\det\Sigma_h(t)}}\,dt,\qquad B_3=\frac{3\,\Gamma(7/6)}{2^{4/3}\pi^{7/2}} .
\end{equation}
The formula is independent of the transverse-frame choice.

Moreover,
\begin{equation}\tag{2.10}\label{eq:correction}
\frac{c_{3,24}}{c_{3,\infty}}-1=-\frac{620813376}{35}\,e^{-288}+\varepsilon_{24},\qquad |\varepsilon_{24}|<10^{-216},
\end{equation}
where $c_{3,\infty}$, the value of \eqref{eq:c} for the Bargmann--Fock field on $\R^3$, is
\begin{equation}\tag{2.11}\label{eq:cinf}
\begin{aligned}
c_{3,\infty}&=2^{-23/6}3^{-8/3}(29\sqrt6-36)\,\Gamma(1/6)\,\pi^{-5/2}\\
&=0.0417759318405983433429366654285755564666815196\ldots .
\end{aligned}
\end{equation}
\end{theorem}

In particular $c_{3,24}<c_{3,\infty}$:
\[
\frac{c_{3,24}}{c_{3,\infty}}-1=-1.48621702906582760722\ldots\cdot10^{-118}+\varepsilon_{24},
\]
so the displayed correction is the exact leading term to a relative precision better than $10^{-98}$.

The leading term of \eqref{eq:main}, with the formula \eqref{eq:c}, is proved in \S\S\ref{ssec:pairs}--\ref{ssec:coefficient}; the value \eqref{eq:cinf} in \S\ref{ssec:flat}; the correction \eqref{eq:correction} in \S\ref{ssec:correction}; and the rate $O(\ell^{1/3})$ in \S\ref{ssec:rate}.

\begin{convention}[density]\label{conv:density}
Throughout, ``density'' means the $\ell$-density of the expected counting measure of finite nonessential superlevel $H_0$ bars per unit volume, with the Kac--Rice counting measure normalized as in Proposition~\ref{prop:density} and \S\S\ref{ssec:pairs}--\ref{ssec:pushforward} (marked two-point formula, ordered typed pairs, per-volume normalization $24^{-3}\int dx=1$). Pointwise statements about $\nu_{3,24}$ and about the candidate and rejected densities of Remark~\ref{rem:cand}, here and below, refer to their canonical marked Kac--Rice density versions: the representatives given by the marked two-point formula, written in the pair coordinates of \S\ref{ssec:pairs} for separations below a fixed $\varrho\in(0,12)$ and in the coordinates $(M,S)$ for the larger separations (Supplement (S6.7) and (S6.10) for $\nu_{3,24}$, and the same formulas without the elder mark for the candidate density). The substitution $\kappa=6\ell/r^3$ shows that their sum does not depend on $\varrho$. A change of density on a null set need not preserve a pointwise asymptotic. All $o(\cdot)$ and $O(\cdot)$ statements are as $\ell\downarrow0$, uniform on compact positive mark sets; the statements about the densities themselves carry no marks.
\end{convention}

\subsection{Beyond the leading term}

The rate in \eqref{eq:main} is not sharp. Theorem~\ref{thm:three-term} is proved in the records listed in Appendix~\ref{app:records}; Part~II of the Supplement states the results it uses, with proofs given in full or condensed. Remarks~\ref{rem:beyond} and~\ref{rem:cand} report further results of those records, at the review status and scope stated with each. None of them is used in the proof of Theorem~\ref{thm:main}; their review status is described in \S\ref{sec:disclosure}.

\begin{theorem}[three-term law]\label{thm:three-term}
For the canonical marked Kac--Rice density version of $\nu_{3,24}$,
\begin{equation}\tag{2.12}\label{eq:three-term}
\nu_{3,24}(\ell)=c_{3,24}\,\ell^{-1/3}+c_1\,\ell^{1/4}+c_2\,\ell^{1/3}+O(\ell^{3/7}),\qquad \ell\downarrow0,
\end{equation}
where the constants are as follows.
\begin{enumerate}[(i)]
\item $c_1<0$ is the cusp coefficient
\begin{equation}\tag{2.13}\label{eq:c1}
c_1=-\frac{192}{7}\,2^{1/4}\int_{S^2}\int_{\R}\pi_0\bigl(u;v_0(b,0)\bigr)\,
\E_0\Bigl[|Y|^{7/4}\,|\det A|^{1/4}\,\ind_{\{A<0\}}\ \Big|\ b\Bigr]\,db\,d\sigma(u),
\end{equation}
where, at a zero-gap contact in direction $u$, $A=D_\perp^2f$ is the transverse Hessian, $\gamma=\nabla_\perp\partial_u^2f$, $f_4=\partial_u^4f$ and $Y=\tfrac{f_4}{12}\det A-\tfrac14\gamma^{\top}\adj(A)\gamma$, and $\pi_0$, $\E_0$ are the zero-gap pin density and conditional expectation of \rec{chain} (Theorem~CU).
\item $c_2$ is the fold-scale finite part of the $r^2$-coefficient of the two-point kernel. At fixed birth height $b$, gap variable $k$ (the gap is $kr^3$, so $k=\kappa/6$ in the notation of \eqref{eq:pins}) and direction $u$, the two-point kernel $\mathcal A_r(b,k,u)$ of \rec{chain} has the fold expansion $\mathcal A_r=\mathcal A_0+r^2\mathcal A_2+O(r^3(1+k^{-1}))$, and
\[
c_2=\frac13\int_{S^2}\int_{\R}\int_0^\infty\bigl[\mathcal A_2(b,k,u)-\mathcal A_2(b,0,u)\bigr]\,k^{-4/3}\,dk\,db\,d\sigma(u),
\]
the Hadamard finite part at $k=0$ of $\tfrac13\int\mathcal A_2\,k^{-4/3}$ \rec{chain}. In the notation of \S\ref{sec:proof}, the kernel is $\mathcal A_r(b,\kappa/6,t)=G^{\mathrm{all}}_r(t,b,\kappa)$, by Supplement~\S S6.
\item Enclosures at $L=24$. First, $c_1/c_{3,24}\in[-5.07106215,\,-5.07106213]$, a certified enclosure of the integral \eqref{eq:c1} divided by $c_{3,24}$, given the enclosure $0.04177593184059834334<c_{3,24}<0.04177593184059834335$ of \rec{v1coef} \rec{consts}; Supplement Corollary S9.16 is that enclosure, and its reading as the coefficient of $\ell^{7/12}$ in \eqref{eq:relative} uses Theorem~CU (Supplement Theorem S9.2 and Remark S9.17). Second, the fixed-cone coefficient expression $c_2[K_{\rm ref}]$ of the Gaussian kernel $K_{\rm ref}(z)=e^{-|z|^2/2}$ on $\R^3$ (Supplement Definition S11.2; \rec{consts}) has a closed form with a certified enclosure,
\[
c_2[K_{\rm ref}]\in[0.1612340491269447,\ 0.1612340491269810]
\]
(Supplement Theorem S11.5). The same expression for the side-24 covariance, $c_2[K_{24}]$, differs from it by less than $10^{-60}$ (Supplement Corollary S11.10), and $c_2=c_2[K_{24}]$: Supplement Corollary S11.11 derives from the fold expansion of Supplement~\S S10 the two identifications that \rec{consts} states as interfaces. So $c_2\in[0.1612340491269446,\ 0.1612340491269811]$. Third, with the enclosure of $c_{3,24}$ in \rec{v1coef},
\[
c_2/c_{3,24}\in[3.859496174547456,\ 3.859496174548331]
\]
(Supplement Corollary S11.12).
\end{enumerate}
Equivalently,
\begin{equation}\tag{2.14}\label{eq:relative}
\nu_{3,24}(\ell)=c_{3,24}\,\ell^{-1/3}\Bigl(1+\tfrac{c_1}{c_{3,24}}\,\ell^{7/12}+\tfrac{c_2}{c_{3,24}}\,\ell^{2/3}+O(\ell^{16/21})\Bigr),\qquad
\tfrac{c_1}{c_{3,24}}\approx-5.0711,\quad \tfrac{c_2}{c_{3,24}}\approx3.8595 .
\end{equation}
\end{theorem}

The proof combines Theorem~R$^+$ (the bounded remainder of Theorem~\ref{thm:main} vanishes for elder pairs), Theorem~R$^{++}$ (its rate $\ell^{1/4}$), Theorem~CU (the cusp term), Theorem~T and Theorems~E3, E3$^+$ (the fold-scale term and the remainder $\ell^{3/7}$) with the far elder bound of \rec{far}; see \rec{chain}. The mechanism behind $c_1$ and $c_2$ is described in \S\ref{ssec:corrections}.

\begin{remark}[the next order]\label{rem:beyond}
The records \rec{twothirds} carry the expansion one order further.
\begin{enumerate}[(a)]
\item The remainder in \eqref{eq:relative} is $O(\ell)$:
\begin{equation}\tag{2.15}
\nu_{3,24}(\ell)=c_{3,24}\,\ell^{-1/3}\Bigl(1+\tfrac{c_1}{c_{3,24}}\ell^{7/12}+\tfrac{c_2}{c_{3,24}}\ell^{2/3}+O(\ell)\Bigr).
\end{equation}
In particular the elder density has no $\ell^{1/2}$ term.
\item More precisely,
\begin{equation}\tag{2.16}\label{eq:next}
\nu_{3,24}(\ell)=c_{3,24}\,\ell^{-1/3}+c_1\ell^{1/4}+c_2\ell^{1/3}+(c_3-R_{2/3})\,\ell^{2/3}+o(\ell^{2/3}),
\end{equation}
where $c_3$ is the $\ell^{2/3}$ coefficient of the candidate density (Remark~\ref{rem:cand}) and $R_{2/3}$ that of the rejected density; both are explicit integrals.
\item At $L=24$, and at the conditional scope of [\ref{rec:soft}, Theorem~A],
\[
(c_3-R_{2/3})(24)\in[0.0773381,\,0.0775860],\qquad\text{i.e.}\quad 0.07746\pm0.00013\approx1.854\,c_{3,24},
\]
with $c_3(24)=0.016092311\pm10^{-9}$ and $R_{2/3}(24)<0$. These enclosures are computer-assisted, in standard-library interval arithmetic. So the coefficient of $\ell$ in the relative expansion is positive.
\end{enumerate}
\end{remark}

\begin{remark}[candidate, elder and rejected densities]\label{rem:cand}
Call \emph{candidate pairs} all ordered typed (maximum, index-two saddle) pairs at every separation, counted by the two-point Kac--Rice formula with the normalization of Convention~\ref{conv:density}, with density $\nu^{\cand}_{3,24}$, and let $\rho_{\rej}=\nu^{\cand}_{3,24}-\nu_{3,24}$ be the density of the candidate pairs that the elder rule rejects. Without the compact-mark restriction of Corollary~\ref{cor:selection}, $\rho_{\rej}$ is not $o(1)$ (Theorem~U of \rec{unres}): it is bounded (Theorem~K and Corollary~T there) and converges, as $\ell\downarrow0$, to a positive constant $B_{3,24}$, an explicit Kac--Rice integral of the equal-height unmarked two-point kernel (Theorem~Z there). Consequently $\rho_{\rej}$ has inverse-moment threshold, the infimum of the exponents $p$ with $\int_0^1\ell^p\rho(\ell)\,d\ell<\infty$, exactly $-1$, whereas $\nu^{\cand}_{3,24}$ and $\nu_{3,24}$ keep the threshold $-2/3$ of the $\ell^{-1/3}$ law. Theorem~\ref{thm:main} concerns the elder density and is unaffected; an earlier conjecture of an $O(\ell^{2/3})$ unrestricted difference is withdrawn \rec{unres}. More precisely,
\begin{equation}\tag{2.17}\label{eq:cand}
\nu^{\cand}_{3,24}(\ell)=c_{3,24}\,\ell^{-1/3}+B_{3,24}+I^{\cand}\ell^{1/4}+c_2\,\ell^{1/3}+c_3\,\ell^{2/3}+O(\ell^{2/3+\delta}),\qquad \delta=\tfrac1{50},
\end{equation}
with $I^{\cand}=\tfrac{3^{1/4}}{2}c_1$ \recs{chain}{parity}; the term $c_3\ell^{2/3}$ and the rate $\delta$ are from notes C3 and C3-R \rec{twothirds}. In particular the candidate density has no $\ell^{1/2}$ term, by a reflection parity of the pinned two-point problem (Theorem~P of \rec{parity}). Subtracting \eqref{eq:three-term} from \eqref{eq:cand},
\begin{equation}\tag{2.18}
\rho_{\rej}(\ell)=B_{3,24}+\bigl(1-\tfrac{3^{1/4}}2\bigr)|c_1|\,\ell^{1/4}+O(\ell^{3/7}),
\end{equation}
so the rejected density approaches its limit from above, and the fold-scale coefficient $c_2$ is common to the candidate and elder densities; by Remark~\ref{rem:beyond} the next term of $\rho_{\rej}$ is $R_{2/3}\ell^{2/3}$. Numerically, as exploration and not as an enclosure, $B_{3,24}\approx10.2$, about $240\,c_{3,24}$: up to a near-field correction of about $-0.02$ it equals $24^3\int\rho_{\max}(b)\rho_{\mathrm{sad}}(b)\,db$, the overlap of the height densities of maxima and index-two saddles \rec{eqh}. Unmarked small-gap pairs on the side-24 torus are thus overwhelmingly unrelated far-apart pairs, which the elder rule removes.
\end{remark}
